\documentclass[12pt,a4paper]{article}
\usepackage[utf8]{inputenc}
\usepackage[spanish]{babel}
\usepackage{amsmath,amsfonts,amssymb}
\usepackage{graphicx}
\usepackage{float}
\usepackage{booktabs}
\usepackage{array}
\usepackage{multirow}
\usepackage[table,xcdraw]{xcolor}
\usepackage{colortbl}
\usepackage{geometry}
\usepackage{fancyhdr}
\usepackage{hyperref}
\usepackage{subcaption}
\usepackage{listings}
\usepackage{enumitem}
\usepackage{tikz}
\usepackage{pgfplots}
\usepackage{setspace}

\geometry{left=2.5cm,right=2.5cm,top=2.5cm,bottom=2.5cm}
\pagestyle{fancy}
\fancyhf{}
\rhead{\thepage}
\lhead{RSL: IA en Agricultura}

\hypersetup{
    colorlinks=true,
    linkcolor=blue,
    filecolor=magenta,
    urlcolor=cyan,
    citecolor=red,
}

\pgfplotsset{compat=1.18}
\onehalfspacing

\title{\textbf{Aplicaciones de Inteligencia Artificial, Visión por Computador y Machine Learning en Agricultura: Una Revisión Sistemática de Literatura}}

\author{
    Álvaro Alejandro Zarabanda Gutierrrez \\
    \small Universidad Distrital Francisco Jose de Caldas \\
    \small 20251595006 \\
    \small Maestría en Ciencias de la Información y las Comunicaciones \\
    \small \texttt{aazarabandag@udistrital.edu.co}
}

\date{\today}

\begin{document}

\maketitle

\begin{abstract}
El sector agrícola enfrenta desafíos críticos relacionados con la seguridad alimentaria global, sostenibilidad ambiental y eficiencia productiva. La aplicación de tecnologías emergentes como inteligencia artificial (IA), visión por computador y machine learning ofrece soluciones prometedoras para abordar estos retos. Esta revisión sistemática de literatura analiza el estado del arte en aplicaciones de IA agrícola, identificando tendencias, metodologías predominantes y oportunidades de investigación futura. Se analizaron 174 publicaciones académicas obtenidas de Scopus, IEEE Xplore, Web of Science y ScienceDirect, correspondientes al período 2014-2024. La metodología incluyó análisis bibliométrico, modelado de tópicos mediante Latent Dirichlet Allocation (LDA), análisis de redes de co-autoría y enriquecimiento de metadatos mediante APIs. Los resultados revelan ... tópicos principales de investigación: .... Se observa . Esta revisión contribuye al entendimiento del panorama actual de la IA agrícola y proporciona direcciones para futuras investigaciones en este campo multidisciplinario.
\end{abstract}

\textbf{Palabras clave:} Inteligencia artificial, Machine learning, Visión por computador, Agricultura de precisión, Detección de enfermedades, Deep learning, Análisis bibliométrico, Revisión sistemática

\newpage

\section{Introducción}

\subsection{Justificación}

\subsection{Preguntas de Investigación}

\subsection{Descripcion del contenido y plan de trabajo}

\section{Metodología}

\subsection{Protocolo de Revisión}

\subsection{Estrategia de Búsqueda}

\subsection{Tesauro de Conceptos}

\subsection{Ecuacion de Búsqueda}

\subsection{Proceso de filtrado de articulos}

\textbf{Criterios de Inclusión}:
\begin{itemize}
    \item Artículos que describan aplicaciones de IA, machine learning o visión por computador en agricultura
    \item Publicaciones en revistas académicas revisadas por pares
    \item Artículos en idioma inglés o español
    \item Estudios que reporten metodologías técnicas específicas y resultados experimentales
    \item Investigaciones que aborden problemas agrícolas reales con soluciones tecnológicas
\end{itemize}

\textbf{Criterios de Exclusión}:
\begin{itemize}
    \item Artículos de revisión, capítulos de libro y publicaciones no revisadas por pares
    \item Estudios puramente teóricos sin validación experimental
    \item Investigaciones que no especifiquen claramente las técnicas de IA utilizadas
    \item Publicaciones duplicadas entre bases de datos
    \item Estudios fuera del dominio agrícola (forestería, acuicultura, etc.)
\end{itemize}

\subsection{Extraccion de datos y sintesis}

\subsection{Proceso de Selección y Extracción de Datos}

El proceso de selección se ejecutó en tres fases: (1) eliminación de duplicados, (2) revisión de títulos y resúmenes, y (3) análisis de texto completo. Se desarrolló un protocolo de extracción de datos que incluyó: información bibliográfica, objetivos del estudio, metodologías técnicas, datasets utilizados, métricas de evaluación, y principales hallazgos.

\section{Resultados y Discusión}

El análisis comprendió múltiples enfoques complementarios:

\textbf{Análisis Bibliométrico}: Se examinaron tendencias temporales de publicación, distribución geográfica, revistas más productivas, y redes de co-autoría utilizando herramientas especializadas.

\textbf{Modelado de Tópicos}: Se aplicó Latent Dirichlet Allocation (LDA) sobre los resúmenes de las publicaciones para identificar temas de investigación emergentes y su evolución temporal.

\textbf{Análisis de Contenido}: Se categorizaron las aplicaciones según el tipo de problema agrícola abordado, las técnicas de IA empleadas, y los tipos de datos utilizados (imágenes, sensores, satelitales, etc.).

\section{Resultados y Discusión}

\subsection{Características Generales de la Literatura}


\subsection{Tópicos de Investigación Identificados}

\subsection{Metodologías Técnicas Predominantes}

\subsection{Tipos de Datos y Plataformas de Adquisición}

\subsection{Desafíos Técnicos y Limitaciones}

\section{Limitaciones del Estudio}

\section{Conclusiones}

\bibliographystyle{plain}

\bibliography{referencias}

\end{document}